% Reader question: Which material must enter the model context now?
% Claim: Staged loading separates discovery, instructions and execution.
% Evidence: Four artifact roles, with conditional retrieval.
% Distinguish: Discovery metadata, task text, references and host code.
% Grammar: A retrieval dependency graph, not invented frequency measurements.
% Acceptance: Identify which text is loaded only when needed.
\pdfigurehue{pdcobalt}%
\begin{tikzpicture}[pd figure,x=1cm,y=1cm]

  \node[pd artifact,text width=6.0cm,align=center] (meta) at (5.2,0) {Discovery metadata\\Name, purpose and activation conditions};
  \node[pd focus state,text width=6.0cm,align=center] (skill) at (5.2,-1.9) {Selected skill instructions\\Load when the task needs them};
  \node[pd artifact,text width=3.7cm,align=center] (refs) at (2.4,-3.9) {Referenced material\\Retrieve relevant sections};
  \node[pd state,text width=3.7cm,align=center] (code) at (8.0,-3.9) {Host-executed scripts\\Return bounded results};
  \draw[pd arrow] (meta) -- (skill);
  \draw[pd arrow] (skill) -- (refs);
  \draw[pd arrow] (skill) -- (code);
  \node[pd note,text width=9.4cm,align=left] at (5.2,-5.5) {Context savings depend on what is retrieved and what the tools return. This architecture alone does not establish a usage frequency or an exponential reduction.};

\end{tikzpicture}
